\chapter{Conclusion}

The proliferation of synthetically generated misinformation represents a pivotal challenge for digital ecosystems. Project Satya successfully demonstrates a functional, GenAI-assisted framework capable of addressing this threat through advanced multimodal verification. Throughout the Software Development Life Cycle, the system transitioned from a conceptual single-script classifier to an asynchronous, microservices-driven architecture operating over text, images, and video.

By systematically applying Google Gemini for nuanced evidence synthesis—and relying on a robust backend infrastructure composed of FastAPI, RabbitMQ, and Redis—the application overcomes the traditional limitations of black-box prediction. Instead, Project Satya actively researches extracted claims, synthesizes context, and presents a structured reasoning report to the user. This \textit{explainability} distinguishes the resulting system not as an infallible truth engine, but as a crucial transparency tool that empowers journalists, researchers, and everyday citizens.

While existing limitations regarding real-time constraints and hallucination risks underscore the continued need for human oversight, the foundation established here is highly extensible. The project achieved its core objective: delivering an accessible, transparent, and scalable system that accurately identifies and justifies verdicts on complex, deceptive media. Ultimately, Project Satya contributes a profound step forward in leveraging artificial intelligence to protect the integrity of human information.
